\chapter{Heat Transfer model}

The heat transfer model is a specific implementation of the \code{Model} interface
dedicated to handle the dynamic heat equation. 
\section{Theory}
The strong form of the heat equation 
can be expressed as
\begin{equation}
  \rho c_v \dot{T} + \nabla \cdot \vec{\kappa} \nabla T = b
\end{equation}
with $T$ the temperature scalar field,
$c_v$ is the specific heat capacity, 
$\rho$ is the mass density, 
$\mat{\kappa}$ is the conductivity tensor and $b$ is the heat generation per unit of volume. 
The weak form discretized with finite elements becomes
\begin{equation}
  \forall i \quad 
  \sum_j \left( \int_\Omega \rho c_v N_j N_i  d\Omega \right) \dot{T}_j 
  - \sum_j \left( \int_\Omega \vec{\kappa} \nabla N_j \nabla N_i d\Omega \right) T_j = 
  - \int_{\Gamma}  N_i \vec{q} \cdot \vec{n} d\Gamma + \int_\Omega b N_i d\Omega
\end{equation}
with $i$ and $j$ being node indexes, $\vec{n}$ the field of normals of the surface 
$\Gamma = \partial \Omega$. 
To simplify, we can define the capacity and conductivity matrices
\begin{equation}
  C_{ij} = \int_\Omega \rho c_v N_j N_i  d\Omega \qquad \textrm{and} \qquad   
  K_{ij} = - \int_\Omega \vec{\kappa} \nabla N_j \nabla N_i d\Omega
\end{equation}
so that the discrete system to solve can be written
\begin{equation}
  \mat{C} \cdot \vec{\dot{T}} = \vec{Q^{ext}} -\mat{K} \cdot \vec{T}~, 
\end{equation}
with $Q^{ext}_i$ the consistent heat generated on node $i$.
\section{Using the heat transfer model}
At the present time an explicit dynamic analysis is available by
using the \code{HeatTransferModel}  object.
It can simply be created like this:
\begin{cpp}
  HeatTransferModel model(mesh, spatial_dimension);
\end{cpp}
while an existing mesh has been used (see \ref{sect:common:mesh}).
Then the model object can be initialized with:
\begin{cpp}
  model.initFull("material.dat")
\end{cpp}
where a material file name has to be provided. This function will load the material 
properties, and allocate/initialize the node and element vectors.
More precisely, the heat transfer model contains 4 \code{Vectors}:
\begin{description}
\item[temperature] contains the nodal temperatures $T$ (zero  by   default  after  the
  initialization) 
\item[temperature\_rate] contains the variation of temperature $\dot{T}$ 
  (zero  by   default  after  the
  initialization) 
%\item[boundary]  contains a  \code{bool} value for  each  node
%  specifying  whether  temperature is  fixed  or  not. A  Dirichlet  boundary
%  condition can be prescribed by setting the \textbf{boundary} value of a node 
%  to  \code{true}.  A Neumann boundary condition  can only be applied
%  if the  \textbf{boundary} value of a node is  \code{false}. If a
%  node has a  \textbf{boundary} value  that is  \code{false}, the
%  \textbf{temperature},     \textbf{temperature\_rate} and
%  \textbf{residual} are computed by the solve algorithm when relevant\todo{wait for till correction}, 
%  otherwise these  vectors  contain   the  imposed  values.

\item[boundary] contains a  boolean value for each degree  of freedom specifying
  whether that degree  is blocked or not. A Dirichlet  boundary condition can be
  prescribed by  setting the \textbf{boundary} value  of a degree  of freedom to
  \code{true}.   A Neumann  boundary condition  can  be applied  by setting  the
  corresponding     \textbf{boundary}     value     to    \code{false}.      The
  \textbf{temperature} and  the \textbf{temperature\_rate} are  computed for all
  degrees  of  freedom   for  which  the  \textbf{boundary}  value   is  set  to
  \code{false}.  For the remaining degrees  of freedom, the imposed values (zero
  by default after  initialization) are kept.

\item[residual] contains the difference between external and internal heat generation. On
  temperature imposed nodes,  \textbf{residual} contains the support heat reactions. ($\vec{R} = \vec{Q^{ext}} -\mat{K} \cdot \vec{T}$)
\end{description}

Only a single material can be specified on the domain. 
A material file allows to give the material properties in a text file (\eg material.dat) as follows:
\begin{cpp}
  heat %\emph{name\_material}% [
  capacity = %\emph{XXX}%
  density = %\emph{XXX}%
  conductivity = [%\emph{XXX}% ... %\emph{XXX}%]
  ]
\end{cpp}
where \code{capacity} and \code{density} are scalars, and the \code{conductivity} is specified as a 
full tensor with $3\times 3$ components.

\subsection{Explicit dynamic}

The implemented explicit  time integration scheme in \akantu  uses a lumped capacity
matrix $\mat{C}$ (reducing the computational  cost, see Chapter \ref{sect:smm}). 
This matrix is assembled by
distributing the capacity of each element onto its nodes. The resulting $\mat{C}$ is
therefore a diagonal matrix stored in the \code{capacity} vector of the model.

\begin{cpp}
  model.assembleCapacityLumped();
\end{cpp}
\index{HeatTransferModel!assembleCapacityLumped}

\note{At current time only explicit time integration with lumped capacity matrix 
is implemented within \akantu.} \\

The explicit integration scheme is a \index{Forward Euler} {\it Forward Euler} 
\cite{curnier92a}.

\begin{itemize}
\item Predictor: $\vec{T}_{n+1} = \vec{T}_{n} + \Delta t \dot{\vec{T}}_{n}$
\item Update residual: $\vec{R}_{n+1} = \left( \vec{Q^{ext}_{n+1}} - \vec{K}\vec{T}_{n+1} \right)$
\item Corrector : $\dot{\vec{T}}_{n+1} = \mat{C}^{-1} \vec{R}_{n+1}$
\end{itemize}

The explicit integration scheme is conditionally stable. The time step has to be
smaller than the stable time step which is obtained in \akantu as follows:

\begin{cpp}
  time_step = model.getStableTimeStep();
\end{cpp}
\index{HeatTransferModel!StableTimeStep}

The stable time step is defined as:
\begin{equation}\label{eqn:smm:explicit:stabletime}
  \Delta t_{\st{crit}} = 2 \Delta x^2 \frac{\rho c_v}{\mid\mid \mat{\kappa} \mid\mid^\infty}
\end{equation}
where $\Delta  x$ is a  characteristic length (\eg  the inradius in the  case of
linear triangle  element), $\rho$ is the density, $\mat{\kappa}$ is the conductivity tensor
and $c_v$ is the specific heat capacity. It is
necessary to impose a  time step that is smaller than the  stable time step, for
instance, by  multiplying the stable time  step by a safety  factor smaller than
one.

\begin{cpp}
  const Real safety_time_factor = 0.1;
  Real applied_time_step = time_step * safety_time_factor;
  model.setTimeStep(applied_time_step);
\end{cpp}

%The initial temperature  and temperature\_rate fields are, by default,  equal to zero if
%not given specifically by the user (similar to the SolidMechanicsModel initial conditions
%\ref{sect:smm:initial_condition}).

The following loop  allows, for each time  step, to update the  temperatures, residual and
temperature\_rate  fields  following the previously described integration scheme.

\begin{cpp}
  for (UInt s = 1; (s-1)*applied_time_step < total_time; ++s) {
    model.explicitPred();
    model.updateResidual();
    model.explicitCorr();  
  }
\end{cpp}
\index{HeatTransferModel!explicitPred}
\index{HeatTransferModel!explicitCorr}
\index{HeatTransferModel!updateResidual}
\index{HeatTransferModel!solveExplicitLumped}

An    example    of    explicit     dynamic    heat propagation is    presented    in \\
\shellcode{\examplesdir/heat\_transfer/explicit\_heat\_transfer.cc}.  \\
This example  consists of a square 2D plate of \unit{1}{\squaremetre} 
having an initial temperature of \unit{100}{\kelvin} everywhere but a none centered hot point 
maintained at \unit{300}{\kelvin}. Figure \ref{fig:htm:explicit:dynamic} presents the
geometry of this case. The material used is a linear fictitious elastic material
with  a density  of  \unit{8940}{\kilogrampercubicmetre}, a  conductivity of 
\unit{401}{\watt\per\metre\per\kelvin} and a specific heat capacity of \unit{385}{\joule\per\kelvin\per\kilogram}. The time step used is \unit{0.12}{\second}.

\begin{figure}[!htb]
  \centering
  \includegraphics[width=.4\textwidth]{figures/hot-point-1}
  \hfill
  \includegraphics[width=.4\textwidth]{figures/hot-point-2}
  \caption{Initial temperature field (left) and after 15000 time steps = 30 minutes (right). The lines represent iso-surfaces.}
  \label{fig:htm:explicit:dynamic}
\end{figure}

